\documentclass[../main.tex]{subfiles}
\begin{document}

\chapter{キャッシュの復習}
\lessoninfo{
  video={Lesson 12，動画 1--45},
  textbook={H\&P \cite{hennessy2017} §2.1，§B.1--B.2},
  keywords={局所性の原理，時間的局所性，空間的局所性，キャッシュヒット，キャッシュミス，AMAT，ミスペナルティ，ブロックオフセット，タグ，有効ビット，ダイレクトマップキャッシュ，セットアソシアティブキャッシュ，LRU，ダーティビット},
}

第3〜11章では，プロセッサが毎秒完了する命令の数を増やし続けてきた．
しかし，どの命令もメモリからフェッチされなければならず，多くの命令は
さらにデータをロードまたはストアする．そして第1章で見たように，メモリの
速度はプロセッサの速度からますます引き離されていく（メモリウォール）．
この章では，この差を埋める機構であるキャッシュを復習する．キャッシュは
小容量で高速なメモリであり，プログラムは局所性をもってメモリにアクセス
するので，大容量で低速な主記憶を待たずに大半のアクセスに応えられる．
この章では，キャッシュの性能をどのように測るか，キャッシュをどのように
構成し探索するか，ミスの際にどのブロックを追い出すか，書き込みをどのように
扱うかを述べる．仮想記憶（第13章）と第14章の高度なキャッシュ技法は，
この内容の上に直接成り立つ．

\section{局所性の原理}
\label{sec:l12-locality}

キャッシュが機能するのは，プログラムのメモリアクセスがランダムとは
ほど遠いからである．\source{Lesson 12，動画 1--6；H\&P §2.1；\cite{denning2005}．}

\begin{definition}[局所性の原理]
  プログラムが近い将来に行うことは，最近行ったことに近い可能性が高い，
  という観察を\term[きょくしょせいのげんり]{局所性の原理}[locality principle]%
  という．
\end{definition}

メモリ参照において，「近い」には2つの意味がある．プログラムがアドレス~$X$
にアクセスしたところだとする．

\begin{definition}[時間的局所性と空間的局所性]
  同じアドレス~$X$ に近いうちに再びアクセスする可能性が高いとき，
  プログラムは\term[じかんてききょくしょせい]{時間的局所性}[temporal locality]%
  を示す．$X$ の近くのアドレスに近いうちにアクセスする可能性が高いとき，
  プログラムは\term[くうかんてききょくしょせい]{空間的局所性}[spatial locality]%
  を示す．
\end{definition}

どちらの局所性もプログラムの構造から生じる．ループは同じ命令を繰り返し
実行し（時間的局所性），それらの命令は連続したアドレスからフェッチされる
（空間的局所性）．反復のたびに更新される変数には同じアドレスで何度も
アクセスし，順に処理される配列の要素は連続したアドレスに並んでいる．

\begin{example}
  \label{ex:l12-locality}
  次のループは，アドレス R4 のバッファからアドレス R5 のバッファへ R6
  バイトをコピーする．
  \begin{lstlisting}[language={[x86masm]Assembler}, numbers=none]
Loop:  LB   R3, 0(R4)      ; R3 <- 転送元のバイト
       SB   R3, 0(R5)      ; 転送先のバイト <- R3
       ADDI R4, R4, 1      ; 次の転送元バイト
       ADDI R5, R5, 1      ; 次の転送先バイト
       ADDI R6, R6, -1     ; 残りのバイト数を 1 減らす
       BNE  R6, R0, Loop   ; 全バイトについて繰り返す
\end{lstlisting}
  6 つの命令フェッチは両方の局所性をもつ．フェッチは連続したアドレスに
  向かい，毎回の反復で同じ 6 つのアドレスがフェッチされる．データアクセス
  は，各バッファを順にたどる2本の流れをなす．どのバイトもちょうど1回だけ
  読まれるか書かれるので，データには時間的局所性がない．それでもデータには
  空間的局所性がある．アクセスは遠く離れているかもしれない2つのバッファの
  間を行き来するが，各アクセスは同じバッファへの直前のアクセスに隣接して
  いる．空間的局所性が要求するのは\emph{最近}アクセスしたアドレスの近くへの
  アクセスであり，最後にアクセスしたアドレスの近くである必要はない．
\end{example}

\section{キャッシュ}
\label{sec:l12-caches}

局所性は，大容量のメモリへの高速なアクセスをプロセッサに与える方法を
示唆する．\source{Lesson 12，動画 7--9；H\&P §B.1；\cite{wilkes1965}．}
フェッチ，ロード，ストアのたびに主記憶にアクセスするのは遅い．主記憶
全体をプロセッサと同じ速さにすることもできない．高速なメモリはビット
当たりのコストが高く，大容量のメモリほど本質的にアクセスが遅いからである．
残る選択肢は，少量の高速なメモリをプロセッサ自身の中に置き，最近アクセス
されたデータを，その周辺のデータとともにそこに置くことである．時間的
局所性と空間的局所性により，大半のアクセスはそこでデータを見つける．

\begin{definition}[キャッシュ]
  \term[きゃっしゅ]{キャッシュ}[cache]とは，より大容量で低速なメモリの
  うち最近アクセスされた部分の写しを保持し，大半のアクセスをその写しで
  処理できるようにする，小容量で高速なメモリである．
\end{definition}

プロセッサのすべてのメモリアクセスは，まずキャッシュ内でデータを探す
（\cref{fig:l12-cache-position}）．データがあれば，そのアクセスは
\term[きゃっしゅひっと]{キャッシュヒット}[cache hit]であり，データは
キャッシュ内で読み書きされる．

そうでなければ，そのアクセスは\term[きゃっしゅみす]{キャッシュミス}[cache miss]%
である．データを主記憶から取得してキャッシュにコピーし，そのデータへの
以降のアクセスと，空間的局所性によりその近傍へのアクセスがヒットするように
する．キャッシュはすべてのアクセスで参照されるので，探索そのものが非常に
高速でなければならない．これが，コストに加えてキャッシュを小さくしなければ
ならない2つ目の理由である．小さなキャッシュにすべてを保持することは
できないので，ミスは避けられない．

\begin{figure}[htbp]
  \centering
  % Processor, cache and main memory: every access looks in the cache first;
% only misses go on to main memory.
\begin{tikzpicture}[x=1cm,y=1cm, font=\footnotesize\sffamily, >={Stealth[length=1.8mm]},
  box/.style={draw, align=center, inner sep=3pt}]
  \node[box, fill=ln-box, minimum width=2.0cm, minimum height=1.2cm] (cpu) at (0,0)
    {\iflangja{プロセッサ}{processor}};
  \node[box, fill=ln-accent!15, minimum width=2.0cm, minimum height=1.0cm] (cache) at (4.0,0)
    {\iflangja{キャッシュ}{cache}\\[-1pt]{\scriptsize\iflangja{小容量・高速}{small, fast}}};
  \node[box, fill=white, minimum width=2.6cm, minimum height=1.8cm] (mem) at (8.2,0)
    {\iflangja{主記憶}{main memory}\\[-1pt]{\scriptsize\iflangja{大容量・低速}{large, slow}}};
  \draw[<->] (cpu) -- node[above, font=\sffamily\scriptsize] {\iflangja{すべてのアクセス}{every access}}
    node[below, font=\sffamily\scriptsize, text=ln-muted] {\iflangja{ヒット時間}{hit time}} (cache);
  \draw[<->, dashed] (cache) -- node[above, font=\sffamily\scriptsize] {\iflangja{ミス時のみ}{misses only}}
    node[below, font=\sffamily\scriptsize, text=ln-muted] {\iflangja{ミスペナルティ}{miss penalty}} (mem);
\end{tikzpicture}

  \caption{すべてのメモリアクセスはキャッシュに向かい，ミスだけが主記憶へ
    進む．}
  \label{fig:l12-cache-position}
\end{figure}

\section{キャッシュの性能}
\label{sec:l12-performance}

キャッシュの効果は，メモリアクセス1回にかかる平均時間で測る．
\source{Lesson 12，動画 10--12；H\&P §B.2．}この時間は次の量で決まる．

\begin{definition}[ヒット率，ミス率]
  \term[ひっとりつ]{ヒット率}[hit rate]はメモリアクセスのうちキャッシュ
  ヒットである割合，\term[みすりつ]{ミス率}[miss rate]はミスである割合で
  ある．ヒット率 $+$ ミス率 $= 1$ である．
\end{definition}

\begin{definition}[ヒット時間，ミスペナルティ]
  \term[ひっとじかん]{ヒット時間}[hit time]とは，キャッシュを探索して
  その中のデータにアクセスするのに必要な時間である．
  \term[みすぺなるてぃ]{ミスペナルティ}[miss penalty]とは，ミスが
  \emph{追加で}要する時間，すなわち主記憶からデータを取り出してキャッシュに
  格納する時間である．
\end{definition}

ミスを検出する前にキャッシュを探索しなければならないので，ミスもヒット
時間を払う．したがってミスにはヒット時間とミスペナルティの和がかかる．
2つの場合をその頻度で重み付けすると，
\term[へいきんめもりあくせすじかん]{平均メモリアクセス時間}[average memory access time]%
\index{AMAT|see{平均メモリアクセス時間}}（AMAT）が得られる．
\begin{equation}
  \begin{aligned}
    \text{AMAT}
    &= (1-\text{ミス率}) \times \text{ヒット時間}
       + \text{ミス率} \times (\text{ヒット時間} + \text{ミスペナルティ}) \\
    &= \text{ヒット時間} + \text{ミス率} \times \text{ミスペナルティ} .
  \end{aligned}
  \label{eq:l12-amat}
\end{equation}
\caution{ミスペナルティはミスの\emph{余分な}時間である．ミスの総時間から
ヒット時間を引くこと．}

AMAT を小さくするには3つのことが必要である．
\begin{itemize}
  \item ヒット時間が短いこと．キャッシュは小さく単純でなければならない．
  \item ミス率が低いこと．キャッシュは，プログラムがこれから使うデータを
    保持できるだけ十分に大きく，かつ巧みに構成されていなければならない．
  \item ミスペナルティが小さいこと．ミスペナルティは主記憶のアクセス時間に
    支配され，大きい．第14章では，キャッシュのレベルを増やすことでこれを
    減らす．
\end{itemize}
最初の2つの目標は相反する．大きなキャッシュはミスが少ないが，探索に時間が
かかる．最適なキャッシュサイズは両者のバランスをとったものであり，その
バランスを測る尺度が AMAT である．

\begin{example}
  \label{ex:l12-amat}
  主記憶によるミスペナルティは \SI{60}{\nano\second} である．3つの候補
  キャッシュを比較する．
  \begin{center}\small
    \begin{tabular}{@{}lccc@{}}
      \toprule
      キャッシュ & ヒット時間 & ミス率 & AMAT \\
      \midrule
      小 & \SI{1}{\nano\second}   & \SI{6}{\percent} & $1 + 0.06\times 60 = \SI{4.6}{\nano\second}$ \\
      中 & \SI{1.5}{\nano\second} & \SI{3}{\percent} & $1.5 + 0.03\times 60 = \SI{3.3}{\nano\second}$ \\
      大 & \SI{3}{\nano\second}   & \SI{2}{\percent} & $3 + 0.02\times 60 = \SI{4.2}{\nano\second}$ \\
      \bottomrule
    \end{tabular}
  \end{center}
  中程度のキャッシュが最良である．大きなキャッシュはミスが少ないが，
  ヒット時間が長いことによるコストが，減らせるミスの分を上回る．3つとも，
  毎回主記憶にアクセスする場合より1桁優れている．
\end{example}

\begin{note}
  キャッシュが引き合うのは，ヒットが主記憶へのアクセスよりずっと速く，
  かつミスが稀な場合だけである．ヒット時間が主記憶のアクセス時間に近ければ，
  ヒットでさえキャッシュなしのアクセスとほぼ同じくらい遅い．ミスが頻繁なら，
  ミスペナルティが支配的になる．どちらの場合も，AMAT は主記憶だけの
  アクセス時間に近づくか，それを上回ることさえある．
\end{note}

\subsection{実際のプロセッサのキャッシュ}

実際のプロセッサはキャッシュを複数のレベルもつ．\source{Lesson 12，動画
13；H\&P §2.1．}最初に参照され，プロセッサの読み書きに直接応えるキャッシュ
が最も小さく高速であり，そのミスだけが，より大きな後続のキャッシュレベル
（L2，L3；第14章），そして最後に主記憶へと進む．この第1レベルを
\term[L1 きゃっしゅ]{L1 キャッシュ}[L1 cache]と呼ぶ．L1 キャッシュは
典型的には 16〜64~KB（1~KB $= 2^{10}$ バイト）であり，アクセスの約
\SI{90}{\percent} がヒットするだけの大きさと，ヒットが 1〜3 クロック
サイクルで済むだけの小ささをもつ．

\section{ブロックとライン}
\label{sec:l12-blocks}

データがキャッシュにあるかどうかを素早く判定するため，キャッシュは固定
サイズのエントリからなる表として構成される．\source{Lesson 12，動画
14--18；H\&P §B.1．}データは主記憶とキャッシュの間を1バイトずつではなく，
ブロックと呼ばれる単位で移動する．

\begin{definition}[キャッシュブロック]
  \term[きゃっしゅぶろっく]{キャッシュブロック}[cache block]とは，
  ひとまとまりとしてキャッシュにコピーされ，またキャッシュから追い出される，
  メモリ上の連続したバイトの並びである．ブロックのバイト数を
  \term[ぶろっくさいず]{ブロックサイズ}[block size]という．
\end{definition}

ブロックサイズは，空間的局所性とキャッシュの容量とのトレードオフである．
\begin{itemize}
  \item ブロックが非常に小さい（極端な場合は 1 バイト）と，空間的局所性
    により起こりやすい次のバイトへのアクセスが再びミスし，バイトごとに
    管理情報が必要になる．
  \item ブロックが非常に大きいと，ブロックの大部分はすぐには使われない
    データになる．それは，最近使われた他のデータを保持できたはずの領域を
    占め，キャッシュに入るブロックの数が減る．
\end{itemize}
典型的なブロックサイズは 32〜128 バイトである．キャッシュが保持できる
ブロックの数は，キャッシュのサイズをブロックサイズで割ったものである．
64 バイトのブロックをもつ 16~KB のキャッシュは $2^{14}/2^{6} = 256$ 個の
ブロックを保持する．

\subsection{ブロックの開始位置}

ブロックは任意のアドレスから始まるわけではない．もしそうなら，アドレス
100〜227 を覆うブロックと 200〜327 を覆うブロックが同時にキャッシュに
存在しうる．すると1つのアドレスが複数のキャッシュ内ブロックに含まれ，
あるアドレスがキャッシュされているかを調べるにはアドレス範囲の比較が
必要になる．そこで，メモリはブロックサイズの倍数から始まる，重なりのない
固定のブロックに分割される．128 バイトのブロックなら，ブロックはアドレス
0，128，256，\ldots\ から始まる（\cref{fig:l12-blocks}）．こうすれば，
どのアドレスもちょうど1つのブロックに属する．

\begin{definition}[キャッシュライン]
  \term[きゃっしゅらいん]{キャッシュライン}[cache line]とは，ブロックを
  1つ保持できるキャッシュ内のスロットである．メモリはブロックに，
  キャッシュはラインに分割される．ラインはブロックと同じ大きさをもち，
  時によって異なるブロックを保持する．
\end{definition}

ブロックサイズ，したがってラインのサイズは，常に2のべき乗である．そうして
初めて，次節で示すように，アドレスが属するブロックとブロック内の位置を，
除算ではなくアドレスのビットの選択によって求められる．

\begin{figure}[htbp]
  \centering
  % Memory is divided into 128-byte blocks starting at multiples of the block
% size; the cache into lines holding one block each, with valid bit and tag.
\begin{tikzpicture}[x=1cm,y=1cm, font=\footnotesize\sffamily, >={Stealth[length=1.8mm]}]
  % --- address fields
  \draw[fill=ln-accent!15] (1.2,0) rectangle (6.4,0.55);
  \node at (3.8,0.275) {\iflangja{ブロック番号}{block number}\ (31\,\dots\,7)};
  \draw[fill=ln-box] (6.4,0) rectangle (9.6,0.55);
  \node at (8.0,0.275) {\iflangja{ブロックオフセット}{block offset}\ (6\,\dots\,0)};
  \node[anchor=east] at (1.15,0.275) {\iflangja{アドレス}{address}};
  % --- main memory
  \node[font=\sffamily\scriptsize, text=ln-muted, anchor=south] at (2.3,-1.0) {\iflangja{主記憶}{main memory}};
  \foreach \i/\a in {0/0x000,1/0x080,2/0x100,3/0x180,4/0x200} {
    \draw[fill=white] (1.2,-1.0-0.5*\i) rectangle (3.4,-1.5-0.5*\i);
    \node[font=\sffamily\scriptsize\ttfamily, anchor=east, text=ln-muted] at (1.15,-1.0-0.5*\i) {\a};
    \node at (2.3,-1.25-0.5*\i) {\iflangja{ブロック}{block}\ \i};
  }
  \fill[ln-accent!15] (1.2,-1.5) rectangle (3.4,-2.0);
  \fill[ln-accent!15] (1.2,-3.0) rectangle (3.4,-3.5);
  \foreach \i in {1,4} {
    \draw (1.2,-1.0-0.5*\i) rectangle (3.4,-1.5-0.5*\i);
    \node at (2.3,-1.25-0.5*\i) {\iflangja{ブロック}{block}\ \i};
  }
  \node at (2.3,-3.75) {$\vdots$};
  % --- cache lines
  \node[font=\sffamily\scriptsize, text=ln-muted, anchor=south] at (6.35,-1.0) {V};
  \node[font=\sffamily\scriptsize, text=ln-muted, anchor=south] at (7.15,-1.0) {\iflangja{タグ}{tag}};
  \node[font=\sffamily\scriptsize, text=ln-muted, anchor=south] at (8.9,-1.0) {\iflangja{データ（128 バイト）}{data (128 bytes)}};
  \foreach \i/\v/\t/\d in {0/1/4/{\iflangja{ブロック 4 の写し}{copy of block 4}},
                           1/0/?/?,
                           2/1/1/{\iflangja{ブロック 1 の写し}{copy of block 1}},
                           3/0/?/?} {
    \draw[fill=white] (6.1,-1.0-0.5*\i) rectangle (6.6,-1.5-0.5*\i);
    \draw[fill=white] (6.6,-1.0-0.5*\i) rectangle (7.7,-1.5-0.5*\i);
    \draw[fill=ln-box] (7.7,-1.0-0.5*\i) rectangle (10.1,-1.5-0.5*\i);
    \node at (6.35,-1.25-0.5*\i) {\v};
    \node at (7.15,-1.25-0.5*\i) {\t};
    \node[font=\sffamily\scriptsize] at (8.9,-1.25-0.5*\i) {\d};
    \node[font=\sffamily\scriptsize, text=ln-muted, anchor=west] at (10.15,-1.25-0.5*\i) {\iflangja{ライン}{line}\ \i};
  }
  \node[font=\sffamily\scriptsize, text=ln-muted, anchor=north] at (8.1,-3.05) {\iflangja{キャッシュ}{cache}};
  % --- copies
  \draw[->, dashed, ln-accent] (3.4,-1.75) .. controls (4.8,-1.75) and (4.6,-2.25) .. (6.1,-2.25);
  \draw[->, dashed, ln-accent] (3.4,-3.25) .. controls (4.9,-3.25) and (4.4,-1.25) .. (6.1,-1.25);
\end{tikzpicture}

  \caption{128 バイトのブロックに分割された主記憶と，4 ラインの
    キャッシュ．ライン~0 と~2 はブロック~4 と~1 の写しを保持し，
    ライン~1 と~3 は有効なものを何も保持していない．}
  \label{fig:l12-blocks}
\end{figure}

\section{ブロック番号，タグ，有効ビット}
\label{sec:l12-tags}

ブロックサイズが $2^{k}$ バイトのとき，アドレス~$A$ は2つのフィールドに
分かれる．\source{Lesson 12，動画 19--23；H\&P §B.1．}

\begin{definition}[ブロックオフセットとブロック番号]
  アドレスの下位 $k$ ビットは\term[ぶろっくおふせっと]{ブロックオフセット}[block offset]%
  であり，そのバイトのブロック内での位置を表す．残りの上位ビットは
  \term[ぶろっくばんごう]{ブロック番号}[block number]であり，ブロックを
  識別する．
  \begin{equation}
    \text{ブロック番号} = \bigl\lfloor A / 2^{k} \bigr\rfloor ,
    \qquad
    \text{ブロックオフセット} = A \bmod 2^{k} .
    \label{eq:l12-block-number}
  \end{equation}
\end{definition}

\begin{example}
  \label{ex:l12-block-number}
  128 バイトのブロック（$k=7$）では，アドレス \texttt{0x0040A7C6} の
  ブロックオフセットは
  $\texttt{0xC6} \mathbin{\&} \texttt{0x7F} = \texttt{0x46} = 70$，
  ブロック番号は $\texttt{0x0040A7C6} \gg 7 = \texttt{0x814F}$ である．
  このバイトはブロック \texttt{0x814F} の位置~70 にあり，このブロックは
  アドレス \texttt{0x0040A780} から \texttt{0x0040A7FF} までを覆う．
\end{example}

\subsection{タグ}

ラインは時によって異なるブロックを保持するので，各ラインは現在どの
ブロックを保持しているかを記録しなければならない．この記録がラインの
\term[たぐ]{タグ}[tag]である．タグは，ラインにあるブロックのブロック番号，
あるいは以降の節で示すように，ブロック番号のうちラインの位置からまだ
分からない部分である．探索では，アドレスのブロック番号を計算し，その
ブロックが置かれうるラインのタグと比較する．一致すればヒットであり，
ブロックオフセットがライン内のバイトを選ぶ．

\subsection{有効ビット}

プロセッサの電源投入時，キャッシュのラインには任意のビット，したがって
任意のタグが入っている．たまたまブロック番号と一致したタグは，一度も
読み込まれていないデータに対するヒットを生んでしまう．

\begin{definition}[有効ビット]
  ラインの\term[ゆうこうびっと]{有効ビット}[valid bit]は，そのラインが
  ブロックを保持しているかどうかを示す．起動時にすべての有効ビットを
  クリアし，ラインにブロックを読み込んだときにそのラインの有効ビットを
  セットする．
\end{definition}

したがって，ライン，あるいはキャッシュ全体を空にするには有効ビットを
クリアするだけでよく，タグとデータはそのままにしておける．唯一の例外は，
キャッシュ内で書き込まれたがまだ主記憶には書き込まれていないブロックで
あり（\cref{sec:l12-write}），これは先に書き戻さなければならない．

探索がヒットとなるのは，ラインが有効で，かつそのタグが一致するときだけ
である．
\begin{equation}
  \text{ヒット} \;=\; V \wedge (\text{タグ} = \text{ブロック番号}) .
  \label{eq:l12-hit}
\end{equation}

\begin{keypoint}
  アドレスはブロック番号とオフセットからなる．タグはラインがどのブロックを
  保持しているかを，有効ビットはそもそもブロックを保持しているかを示し，
  オフセットはバイトを選ぶ．キャッシュ構成におけるそれ以外のすべては，
  \emph{どの}ラインを探索しなければならないかを決めるためのものである．
\end{keypoint}

\section{ブロックを置ける場所}
\label{sec:l12-placement}

キャッシュは，あるブロックを置けるラインがいくつあるかによって
異なる．\source{Lesson 12，動画 24；H\&P §B.1．}この数をキャッシュの
\term[れんそうど]{連想度}[associativity]という．連想度は，探索で比較
しなければならないタグの数と，異なるブロックが同じラインを奪い合う頻度を
決める．次の3つの構成が区別される（\cref{fig:l12-placement}）．
\begin{description}
  \item[ダイレクトマップ] 各ブロックはちょうど1つのラインにしか置けない
    （\cref{sec:l12-dm}）．
  \item[セットアソシアティブ] 各ブロックは，セットをなす $N$ 個のラインの
    どれにでも置ける（\cref{sec:l12-sa}）．
  \item[フルアソシアティブ] 各ブロックはキャッシュのどのラインにでも置ける
    （\cref{sec:l12-fa}）．
\end{description}

\begin{figure}[htbp]
  \centering
  % Where block number 13 may be placed in an 8-line cache that is
% direct-mapped, 2-way set-associative, or fully associative.
\begin{tikzpicture}[x=1cm,y=1cm, font=\footnotesize\sffamily]
  \def\rh{0.34}
  \foreach \cx/\ttl/\how in {
      1.7/{\iflangja{ダイレクトマップ}{direct-mapped}}/{\iflangja{ライン}{line} $13 \bmod 8 = 5$},
      5.5/{\iflangja{2 ウェイセットアソシアティブ}{2-way set-associative}}/{\iflangja{セット}{set} $13 \bmod 4 = 1$},
      9.3/{\iflangja{フルアソシアティブ}{fully associative}}/{\iflangja{任意のライン}{any line}}} {
    \node[font=\sffamily\scriptsize\bfseries, anchor=south, text width=3.4cm, align=center] at (\cx,0.08) {\ttl};
    \node[font=\sffamily\scriptsize, anchor=north] at (\cx,-8*\rh-0.1) {\how};
    \foreach \i in {0,...,7} {
      \node[font=\sffamily\scriptsize, text=ln-muted, anchor=east] at (\cx-0.8,-\rh*\i-\rh/2) {\i};
    }
  }
  % highlighted lines
  \fill[ln-accent!15] (0.9,-5*\rh) rectangle (2.5,-6*\rh);
  \fill[ln-accent!15] (4.7,-2*\rh) rectangle (6.3,-4*\rh);
  \fill[ln-accent!15] (8.5,0) rectangle (10.1,-8*\rh);
  % grids
  \foreach \cx in {1.7,5.5,9.3} {
    \foreach \i in {0,...,7} { \draw (\cx-0.8,-\rh*\i) rectangle (\cx+0.8,-\rh*\i-\rh); }
  }
  % set boundaries and labels in the 2-way cache
  \foreach \s in {0,...,3} {
    \draw[thick] (4.7,-2*\s*\rh) rectangle (6.3,-2*\s*\rh-2*\rh);
    \node[font=\sffamily\scriptsize, text=ln-muted, anchor=west] at (6.35,-2*\s*\rh-\rh) {\iflangja{セット}{set}\ \s};
  }
\end{tikzpicture}

  \caption{各構成の 8 ラインのキャッシュにおいて，ブロック番号~13 を
    置けるライン．}
  \label{fig:l12-placement}
\end{figure}

\section{ダイレクトマップキャッシュ}
\label{sec:l12-dm}

3つの構成のうち最も単純なものは，すべてのブロックにキャッシュのラインを
1つだけ割り当てる．\source{Lesson 12，動画 25--28；H\&P §B.1．}
\term[だいれくとまっぷきゃっしゅ]{ダイレクトマップキャッシュ}[direct-mapped cache]%
では，どのブロックにもそれを置けるラインがちょうど1つあり，そのラインは
ブロック番号の下位ビットで選ばれる．

\begin{definition}[インデックス]
  \term[いんでっくす]{インデックス}[index]とは，ブロックをキャッシュの
  どこに置けるかを選ぶアドレスのビット群である．$2^{n}$ 個のラインをもつ
  ダイレクトマップキャッシュでは，インデックスはブロック番号の下位 $n$
  ビットからなる．
\end{definition}

したがってアドレスは3つのフィールドに分かれる（\cref{fig:l12-direct-mapped}）．
最下位の $k$ ビットがオフセット，その上の $n$ ビットがインデックス，残りの
すべてのビットがタグである．ラインにはブロック番号全体ではなく，このタグの
ビットだけを格納する．ライン~$i$ に置けるブロックはすべてインデックス~$i$
をもつので，ブロック番号のインデックス部分はライン自体から分かるからである．

探索では，インデックスで選ばれた1つのラインを読み，そのタグをアドレスの
タグと比較し，有効ビットを調べる（\cref{eq:l12-hit}）．

\begin{figure}[htbp]
  \centering
  % Direct-mapped lookup for a 64 KB cache with 128-byte blocks and 32-bit
% addresses: the index selects one line, one comparator checks the tag.
\begin{tikzpicture}[x=1cm,y=1cm, font=\footnotesize\sffamily, >={Stealth[length=1.8mm]}]
  % --- address fields
  %% The Japanese index label filled its box edge to edge, so the Japanese
  %% tag/index boundary sits further left (the index arrow at x=3.3 and the
  %% tag arrow at x=1.5 stay inside their boxes).
  \iflangja{%
    \draw[fill=ln-accent!15] (0,0) rectangle (2.6,0.55);
    \node at (1.3,0.275) {タグ\ (31\,\dots\,16)};
    \draw[fill=ln-box] (2.6,0) rectangle (5.6,0.55);
    \node at (4.1,0.275) {インデックス\ (15\,\dots\,7)};
  }{%
    \draw[fill=ln-accent!15] (0,0) rectangle (3.0,0.55);
    \node at (1.5,0.275) {tag\ (31\,\dots\,16)};
    \draw[fill=ln-box] (3.0,0) rectangle (5.6,0.55);
    \node at (4.3,0.275) {index\ (15\,\dots\,7)};
  }
  \draw[fill=white] (5.6,0) rectangle (8.2,0.55);
  \node at (6.9,0.275) {\iflangja{オフセット}{offset}\ (6\,\dots\,0)};
  % --- table: V | tag | data, 6 drawn rows
  \node[font=\sffamily\scriptsize, text=ln-muted, anchor=south] at (3.85,-1.0) {V};
  \node[font=\sffamily\scriptsize, text=ln-muted, anchor=south] at (4.9,-1.0) {\iflangja{タグ}{tag}};
  \node[font=\sffamily\scriptsize, text=ln-muted, anchor=south] at (7.2,-1.0) {\iflangja{データ}{data}};
  \fill[ln-accent!15] (3.6,-2.35) rectangle (8.7,-2.8);
  \foreach \i/\lab in {0/0,1/1,2/{},3/335,4/{},5/511} {
    \draw (3.6,-1.0-0.45*\i) rectangle (4.1,-1.45-0.45*\i);
    \draw (4.1,-1.0-0.45*\i) rectangle (5.7,-1.45-0.45*\i);
    \draw (5.7,-1.0-0.45*\i) rectangle (8.7,-1.45-0.45*\i);
    \node[font=\sffamily\scriptsize, text=ln-muted, anchor=west] at (8.75,-1.225-0.45*\i) {\lab};
  }
  \foreach \i in {2,4} {
    \node[font=\small, text=ln-muted, inner sep=0pt] at (4.9,-1.225-0.45*\i) {$\vdots$};
    \node[font=\small, text=ln-muted, inner sep=0pt] at (7.2,-1.225-0.45*\i) {$\vdots$};
  }
  % --- index selects a line
  \draw[->] (3.3,0) -- (3.3,-2.575) -- (3.6,-2.575);
  % --- tag comparison and valid bit
  \node[draw, circle, fill=white, inner sep=1pt, minimum size=0.5cm] (eq) at (4.9,-4.25) {$=$};
  \node[draw, fill=ln-box, minimum width=0.7cm, minimum height=0.5cm, inner sep=1pt, font=\sffamily\scriptsize] (and) at (3.85,-4.25) {AND};
  \draw[->] (4.9,-3.7) -- (eq.north);
  \draw[->] (3.85,-3.7) -- (and.north);
  \draw[->] (1.5,0) -- (1.5,-5.0) -- (4.9,-5.0) -- (eq.south);
  \draw[->] (eq.west) -- (and.east);
  \draw[->] (and.west) -- (2.6,-4.25) node[left] {\iflangja{ヒット}{hit}};
  % --- data selection by offset
  \node[draw, fill=white, minimum width=2.0cm, minimum height=0.5cm, inner sep=1pt, font=\sffamily\scriptsize] (sel) at (7.2,-4.25)
    {\iflangja{バイトを選択}{select byte}};
  \draw[->] (7.2,-3.7) -- (sel.north);
  \draw[->] (8.2,0.275) -- (9.6,0.275) -- (9.6,-4.25) -- (sel.east);
  \draw[->] (sel.south) -- (7.2,-5.0) node[below] {\iflangja{データ}{data}};
\end{tikzpicture}

  \caption{128 バイトのブロックをもつ 64~KB のダイレクトマップキャッシュに
    おける探索．インデックスが 512 個のラインの1つを選び，1つの比較器が
    そのタグを調べる．}
  \label{fig:l12-direct-mapped}
\end{figure}

\begin{example}
  \label{ex:l12-dm-fields}
  128 バイトのブロックと 32 ビットのアドレスをもつ 64~KB のダイレクト
  マップキャッシュには $512 = 2^{9}$ 個のラインがあり，オフセットは
  ビット 6--0，インデックスはビット 15--7，タグはビット 31--16 である．
  \cref{ex:l12-block-number} のアドレス \texttt{0x0040A7C6}（ブロック番号
  \texttt{0x814F}）のインデックスは
  $\texttt{0x814F} \mathbin{\&} \texttt{0x1FF} = \texttt{0x14F} = 335$，
  タグは $\texttt{0x814F} \gg 9 = \texttt{0x0040}$ である．ライン~335 が
  有効でタグ \texttt{0x0040} を保持していれば，このアクセスはヒットする．
\end{example}

\subsection{利点と欠点}

ダイレクトマップキャッシュの長所は，調べるラインが1つだけで済むことである．
比較器は1つで足りるので，探索は高速かつ安価で，エネルギー効率も良い．
ヒット時間は，与えられたサイズのキャッシュとして可能な限り短い．

短所は，同じインデックスをもつブロックが1つのラインを奪い合うことである．
プログラムがそのような2つのブロックに交互にアクセスすると，各アクセスが
もう一方のアクセスに必要なブロックを追い出し，キャッシュの残りが空で
あってもすべてのアクセスがミスする．こうした競合がミス率を高める．

\begin{example}
  \label{ex:l12-dm-conflict}
  16 バイトのブロックをもつ 8 ラインのダイレクトマップキャッシュでは，
  オフセットはビット 3--0，インデックスはビット 6--4，タグはビット 31--7
  である．あるプログラムがアドレス \texttt{0x1000}（ブロック
  \texttt{0x100}）と \texttt{0x1084}（ブロック \texttt{0x108}）に交互に
  アクセスする．どちらのブロックもインデックスは~0 で，タグはそれぞれ
  \texttt{0x20} と \texttt{0x21} である．各アクセスは直前のアクセスが
  読み込んだブロックを置き換えるので，プログラムがどれだけ長く続いても，
  すべてのアクセスがミスとなる．
\end{example}

\section{セットアソシアティブキャッシュ}
\label{sec:l12-sa}

セットアソシアティブ構成は，ラインが1つだけという制約を
緩める．\source{Lesson 12，動画 29--31；H\&P §B.1．}

\begin{definition}[セットアソシアティブキャッシュ]
  \term[せっとあそしあてぃぶきゃっしゅ]{セットアソシアティブキャッシュ}[set-associative cache]%
  は，ラインを $N$ 個ずつの\emph{セット}に分け，ブロックを特定の1つの
  セットのどのラインにでも置けるようにする．このようなキャッシュを $N$
  ウェイセットアソシアティブという．
\end{definition}

このときインデックスはラインではなくセットを選ぶ．$2^{s}$ 個のセットを
もつキャッシュのインデックスは $s$ ビットであり，タグは再びインデックス
より上のすべてのビットからなる．探索では，アドレスのタグを，選ばれた
セットの $N$ 個すべてのラインのタグと並列に比較し，有効なラインのいずれかが
一致すればヒットとなる（\cref{fig:l12-set-associative}）．データは一致した
ラインから得られる．

\begin{figure}[htbp]
  \centering
  % Lookup in a 2-way set-associative cache with 4 sets (8 lines, 16-byte
% blocks, 32-bit addresses): the index selects a set, both tags are compared,
% and each comparison is ANDed with the valid bit of its line.
\begin{tikzpicture}[x=1cm,y=1cm, font=\footnotesize\sffamily, >={Stealth[length=1.8mm]},
  gate/.style={draw, fill=ln-box, minimum width=0.7cm, minimum height=0.5cm, inner sep=1pt, font=\sffamily\scriptsize},
  cmp/.style={draw, circle, fill=white, inner sep=1pt, minimum size=0.5cm}]
  % --- address fields
  %% The Japanese index label filled its box edge to edge, so the Japanese
  %% tag/index boundary sits further left (the index line at x=5.4 and the
  %% tag line at x=1.0 stay inside their boxes).
  \iflangja{%
    \draw[fill=ln-accent!15] (0.6,0) rectangle (3.6,0.55);
    \node at (2.1,0.275) {タグ\ (31\,\dots\,6)};
    \draw[fill=ln-box] (3.6,0) rectangle (6.4,0.55);
    \node at (5.0,0.275) {インデックス\ (5\,\dots\,4)};
  }{%
    \draw[fill=ln-accent!15] (0.6,0) rectangle (4.0,0.55);
    \node at (2.3,0.275) {tag\ (31\,\dots\,6)};
    \draw[fill=ln-box] (4.0,0) rectangle (6.4,0.55);
    \node at (5.2,0.275) {index\ (5\,\dots\,4)};
  }
  \draw[fill=white] (6.4,0) rectangle (8.8,0.55);
  \node at (7.6,0.275) {\iflangja{オフセット}{offset}\ (3\,\dots\,0)};
  % --- two ways, four sets; set 0 selected
  \foreach \x in {1.6,5.8} {
    \fill[ln-accent!15] (\x,-1.1) rectangle (\x+3.4,-1.55);
    \node[font=\sffamily\scriptsize, text=ln-muted, anchor=south] at (\x+0.2,-1.1) {V};
    \node[font=\sffamily\scriptsize, text=ln-muted, anchor=south] at (\x+1.1,-1.1) {\iflangja{タグ}{tag}};
    \node[font=\sffamily\scriptsize, text=ln-muted, anchor=south] at (\x+2.6,-1.1) {\iflangja{データ}{data}};
    \foreach \i in {0,...,3} {
      \draw (\x,-1.1-0.45*\i) rectangle (\x+0.4,-1.55-0.45*\i);
      \draw (\x+0.4,-1.1-0.45*\i) rectangle (\x+1.8,-1.55-0.45*\i);
      \draw (\x+1.8,-1.1-0.45*\i) rectangle (\x+3.4,-1.55-0.45*\i);
    }
  }
  \node[font=\sffamily\scriptsize, text=ln-muted, anchor=north] at (4.2,-2.95) {\iflangja{ウェイ 0}{way 0}};
  \node[font=\sffamily\scriptsize, text=ln-muted, anchor=north] at (8.4,-2.95) {\iflangja{ウェイ 1}{way 1}};
  \foreach \i in {0,...,3} {
    \node[font=\sffamily\scriptsize, text=ln-muted, anchor=west] at (9.25,-1.325-0.45*\i) {\iflangja{セット}{set}\ \i};
  }
  % --- index selects a set in both ways
  \draw (5.4,0) -- (5.4,-1.325);
  \fill (5.4,-1.325) circle (1.2pt);
  \draw[->] (5.4,-1.325) -- (5.0,-1.325);
  \draw[->] (5.4,-1.325) -- (5.8,-1.325);
  % --- per way: tag comparator, AND with the valid bit
  \foreach \x/\w in {1.6/0,5.8/1} {
    \node[cmp] (c\w) at (\x+1.2,-3.6) {$=$};
    \node[gate] (a\w) at (\x+0.4,-4.35) {AND};
    \draw[->] (\x+1.2,-2.9) -- (c\w.north);
    \draw[->] (\x+0.2,-2.9) -- (\x+0.2,-4.1);
    \draw[->] (c\w.west) -| (\x+0.6,-4.1);
  }
  % --- tag of the address to both comparators
  \draw (1.0,0) -- (1.0,-5.5) -- (7.0,-5.5);
  \draw[->] (7.0,-5.5) -- (c1.south);
  \fill (2.8,-5.5) circle (1.2pt);
  \draw[->] (2.8,-5.5) -- (c0.south);
  % --- either way hits
  \node[gate] (or) at (4.5,-4.35) {OR};
  \draw[->] (a0.east) -- (or.west);
  \draw[->] (a1.west) -- (or.east);
  \draw[->] (or.south) -- (4.5,-4.95) node[below] {\iflangja{ヒット}{hit}};
\end{tikzpicture}

  \caption{4 セットの 2 ウェイセットアソシアティブキャッシュ
    （\cref{ex:l12-sa-conflict} の構成）における探索．インデックス
    （ここではその例と同じ~0）がセットを選び，セット内の2つのラインは
    それぞれ，有効でタグが一致すればヒットとなる．}
  \label{fig:l12-set-associative}
\end{figure}

\begin{example}
  \label{ex:l12-sa-conflict}
  \cref{ex:l12-dm-conflict} の 8 ラインを 2 ウェイセットアソシアティブ
  キャッシュとして構成する．セットは 4 個なので，インデックスはビット
  5--4，タグはビット 31--6 である．ブロック \texttt{0x100} はタグ
  \texttt{0x40} でセット~0 に，ブロック \texttt{0x108} はタグ
  \texttt{0x42} でセット~0 に入る．2つのブロックは依然として同じセットを
  共有するが，セットには2つのラインがある．各ブロックへの最初のアクセスが
  ミスしてそのブロックを読み込んだ後は，以降のアクセスはすべてヒットする．
\end{example}

同じサイズのダイレクトマップキャッシュと比べると，$N$ ウェイセット
アソシアティブキャッシュには $N$ 個の比較器と，一致したラインのデータの
選択が必要であり，そのためヒット時間は長くなる．しかし競合が少ないので，
ミス率は低い．

\section{フルアソシアティブキャッシュ}
\label{sec:l12-fa}

もう一方の極端では，ブロックをどこにでも置ける．\source{Lesson 12，動画
32--33；H\&P §B.1．}

\begin{definition}[フルアソシアティブキャッシュ]
  \term[ふるあそしあてぃぶきゃっしゅ]{フルアソシアティブキャッシュ}[fully associative cache]%
  では，どのブロックもどのラインにでも置ける．
\end{definition}

選ぶものがないので，インデックスは存在しない．タグはブロック番号全体で
あり，探索ではそれをキャッシュ内のすべてのラインのタグと比較する．空いて
いるラインがある限り2つのブロックが1つのラインを奪い合うことはないが，
ラインごとに比較器が必要になるので，大きなフルアソシアティブキャッシュは
低速で高価になる．

\subsection{3つの構成に共通の枠組み}

ダイレクトマップキャッシュとフルアソシアティブキャッシュは，セット
アソシアティブキャッシュの両極端である．ダイレクトマップキャッシュは，
ラインと同数のセットをもつ 1 ウェイセットアソシアティブキャッシュであり，
フルアソシアティブキャッシュは，すべてのラインを含むただ1つのセットを
もつ．アドレスが $a$ ビットのとき，どの構成についても次が成り立つ．
\begin{equation}
  \begin{aligned}
    \text{オフセットビット数} &= \log_2(\text{ブロックサイズ}), \\
    \text{インデックスビット数}  &= \log_2(\text{セット数})
                        = \log_2\bigl(\text{ライン数}/N\bigr), \\
    \text{タグビット数}    &= a - \text{インデックスビット数} - \text{オフセットビット数} ,
  \end{aligned}
  \label{eq:l12-fields}
\end{equation}
ここで，ダイレクトマップキャッシュでは $N=1$（インデックスビット数
$= \log_2(\text{ライン数})$），フルアソシアティブキャッシュでは
$N = \text{ライン数}$（インデックスは 0 ビット）である．
\cref{ex:l12-dm-fields} の 128 バイトのブロックをもつ 64~KB のキャッシュ
では，タグはダイレクトマップの 16 ビットからフルアソシアティブの 25 ビット
まで増える．その中間の場合は \cref{sec:l12-summary} の
\cref{tab:l12-organizations} に示す．\tip{$C$ バイトのキャッシュのタグは
$a - \log_2(C/N)$ ビットである．$N$ を2倍にするたびに，1 ビットが
インデックスからタグに移る．}

\begin{keypoint}
  連想度はヒット時間とミス率のトレードオフを決める．セット当たりのウェイ数
  が多いほど，毎回のアクセスでのタグ比較は増えるが，同じラインを奪い合う
  ブロックは減る．
\end{keypoint}

\section{キャッシュの置換}
\label{sec:l12-replacement}

ミスの際には，新しいブロックを，それに許されたラインのいずれかに置か
なければならない．\source{Lesson 12，動画 34--36；H\&P §B.1．}それらの
ラインのうち有効でないものがあれば，そのラインを使う．すべてが有効な
ブロックを保持していれば，1つのブロックを追い出さなければならない．
ダイレクトマップキャッシュには選択の余地がない．セットアソシアティブ
キャッシュとフルアソシアティブキャッシュでは，追い出すブロックを
\term[ちかんあるごりずむ]{置換アルゴリズム}[replacement policy]が選ぶ．
\begin{description}
  \item[ランダム] セット内のブロックを無作為に選んで追い出す．単純だが，
    ブロックの使われ方を無視する．
  \item[FIFO]（first in, first out）キャッシュに最も長く存在している
    ブロックを，まだ使われているかどうかにかかわらず追い出す．
  \item[LRU] 最後にアクセスされたのが最も昔であるブロックを追い出す．
\end{description}

\begin{definition}[LRU]
  \term{LRU}[least recently used]（least recently used）方式は，置き換えて
  よいブロックのうち，最も長い間アクセスされていないものを追い出す．
\end{definition}

LRU は時間的局所性を利用する．最近アクセスされたブロックは近いうちに
再びアクセスされる可能性が高いので，最も長い間使われていないブロックが，
必要とされる可能性の最も低いブロックである．

\subsection{カウンタによる LRU の実装}

LRU は，ラインごとのカウンタで実装できる．$N$ ウェイセットアソシアティブ
キャッシュでは，各ラインが $\log_2 N$ ビットのカウンタをもつ．起動時に
各セットのカウンタを任意の順序で値 $0, 1, \ldots, N-1$ に初期化し，以下の
更新規則はこの状態を保つ．すなわち，各セットでどの値もちょうど1回ずつ
現れる．カウンタが~0 のラインが最も長い間使われていないラインであり，
カウンタが $N-1$ のラインが最も最近使われたラインである．カウンタが~$c$
であるラインへのアクセスのたびに，次を行う．
\begin{enumerate}
  \item セット内で $c$ より大きいカウンタをすべて 1 減らす．
  \item アクセスしたラインのカウンタを $N-1$ にする．
\end{enumerate}
ミスの際は，カウンタが~0 のラインを追い出す．新しいブロックがそこに
入った後は，そのアクセスを他のアクセスと同様に扱い，上のようにカウンタを
更新する．\cref{tab:l12-lru-trace} は 4 ウェイのセット1つを追ったもので
ある．\caution{カウンタはヒットを含む\emph{すべての}アクセスで更新する
こと．読み込み時にだけ更新すると FIFO になる．}

\begin{table}[htbp]
  \centering\small
  \caption{ブロック A〜D を保持する 4 ウェイのセット1つの LRU カウンタ．
    4つのラインそれぞれについて，ブロック:カウンタの形で示す．}
  \label{tab:l12-lru-trace}
  \begin{tabular}{@{}llcccc@{}}
    \toprule
    アクセス & 結果 & ライン 0 & ライン 1 & ライン 2 & ライン 3 \\
    \midrule
    （初期状態） &                     & A:0 & B:1 & C:2 & D:3 \\
    C         & ヒット              & A:0 & B:1 & C:3 & D:2 \\
    A         & ヒット              & A:3 & B:0 & C:2 & D:1 \\
    E         & ミス，B を追い出す  & A:2 & E:3 & C:1 & D:0 \\
    D         & ヒット              & A:1 & E:2 & C:0 & D:3 \\
    B         & ミス，C を追い出す  & A:0 & E:1 & B:3 & D:2 \\
    \bottomrule
  \end{tabular}
\end{table}

この実装のコストは，ラインごとの $\log_2 N$ ビットと，より重要なことに，
すべてのアクセスで生じる動作である．ヒットでさえセットの $N$ 個のカウンタ
すべてを変更しうる．連想度の高いキャッシュでは，これはかなりのハード
ウェアとエネルギーを要し，第14章の，より安価な LRU の近似の動機となる．

\section{書き込み方針}
\label{sec:l12-write}

書き込みには，読み出しでは生じない問題が2つある．\source{Lesson 12，動画
37--40；H\&P §B.1．}

\begin{enumerate}
  \item \emph{書き込みミスの際にブロックをキャッシュに持ち込むか．}
    持ち込むキャッシュは\term[らいとあろけーとほうしき]{ライトアロケート方式}[write-allocate policy]%
    に従う．ノーライトアロケートのキャッシュは値を主記憶だけに書き込む．
    大半のキャッシュはライトアロケートである．局所性により，書き込まれた
    ばかりのブロックは近いうちに再び読み書きされる可能性が高いからである．
  \item \emph{主記憶が書き込まれた値を受け取るのはいつか．}これには2つの
    方式がある．
\end{enumerate}

\begin{definition}[ライトスルーとライトバック]
  \term[らいとするーほうしき]{ライトスルー方式}[write-through policy]では，
  書き込みはキャッシュ内のブロックを更新し，ただちに主記憶内のブロックも
  更新する．

  \term[らいとばっくほうしき]{ライトバック方式}[write-back policy]では，
  書き込みはキャッシュ内のブロックだけを更新し，主記憶はそのブロックが
  キャッシュから追い出されるときに更新される．
\end{definition}

ライトスルーでは主記憶が常に最新に保たれるが，そのかわりすべての書き込みが
主記憶のアクセス時間を払うことになり，これはまさにキャッシュが避けようと
したものである．そのためキャッシュはライトバックを用いる．ブロックが
キャッシュにある間の多数の書き込みは，追い出し時の高々1回の主記憶への
書き込みで済む．ライトバックのキャッシュはライトアロケートでもある．
そうでなければ，キャッシュにないブロックへの書き込みは，そのブロックが
どれほど頻繁に書き込まれても，すべて主記憶に向かうことになる．

しかし，追い出すブロックをすべて書き戻すと，読み出されただけでメモリ上の
写しと同一のブロックまで書き込むことになる．

\begin{definition}[ダーティビット]
  ラインの\term[だーてぃびっと]{ダーティビット}[dirty bit]は，そのブロック
  が読み込まれてから書き込まれたかどうかを示す．ブロックを読み込んだときに
  クリアし，ブロックが書き込まれたときにセットする．追い出しの際，ダーティ
  なブロックは主記憶に書き込み，クリーンなブロックは単に破棄する．
\end{definition}

\begin{example}
  \label{ex:l12-write-back}
  ブロック X と Y はダイレクトマップキャッシュの同じラインに対応し，この
  ラインは最初は無効である．\cref{tab:l12-write-back} は，ライトバック
  方式のもとでの一連のアクセスを追ったものである．Y への連続した2回の
  書き込みは，ステップ~6 での1回のメモリ書き込みで済み，ステップ~8 で
  Y を追い出してもメモリ書き込みは生じない．Y はステップ~7 で読み込まれた
  後に書き込まれていないからである．ライトバックに必要なメモリ書き込みは
  合計 3 回であり，ライトスルーなら書き込みごとに1回，すなわち 4 回必要
  になる．
\end{example}

\begin{table}[htbp]
  \centering\small
  \caption{ライトバック方式のもとでの1つのライン
    （\cref{ex:l12-write-back}）．D はアクセス後のダーティビットである．}
  \label{tab:l12-write-back}
  \begin{tabular}{@{}cllll@{}}
    \toprule
    ステップ & アクセス  & 結果 & メモリ操作        & アクセス後のライン \\
    \midrule
    1    & X 読み出し & ミス   & X 読み込み                & X，D$=$0 \\
    2    & X 書き込み & ヒット & ---                       & X，D$=$1 \\
    3    & Y 読み出し & ミス   & X 書き戻し，Y 読み込み    & Y，D$=$0 \\
    4    & Y 書き込み & ヒット & ---                       & Y，D$=$1 \\
    5    & Y 書き込み & ヒット & ---                       & Y，D$=$1 \\
    6    & X 書き込み & ミス   & Y 書き戻し，X 読み込み    & X，D$=$1 \\
    7    & Y 読み出し & ミス   & X 書き戻し，Y 読み込み    & Y，D$=$0 \\
    8    & X 読み出し & ミス   & X 読み込み                & X，D$=$0 \\
    \bottomrule
  \end{tabular}
\end{table}

\section{全体のまとめ}
\label{sec:l12-summary}

キャッシュはいくつかの設計上の選択によって完全に記述され，その構造は
そこから導かれる．\source{Lesson 12，動画 41--45；H\&P §B.1．}
\begin{description}
  \item[キャッシュサイズとブロックサイズ] ライン数とオフセットのビット数が
    決まる．
  \item[連想度] セット数，したがってインデックスとタグのビット数が決まる
    （\cref{eq:l12-fields}）．
  \item[置換アルゴリズム] LRU は各ラインに $\log_2 N$ ビットのカウンタを
    加える（ダイレクトマップキャッシュでは不要）．
  \item[書き込み方針] ライトバックは各ラインにダーティビットを加える．
\end{description}
こうして各ラインは，データに加えて有効ビット，ダーティビット，LRU
カウンタ，タグを格納する．\cref{fig:l12-access-flow} は，1回のアクセスで
これらがどのように使われるかを示す．ヒットはデータを読み書きし，LRU
カウンタを更新する．ミスはまず追い出すラインを選び，それが有効かつ
ダーティであれば書き戻してから，新しいブロックを読み込む．その後，
アクセスはヒットと同様に進む．

\begin{figure}[htbp]
  \centering
  % One memory access in a set-associative cache with LRU replacement and
% the write-back policy.
\begin{tikzpicture}[x=1cm,y=1cm, font=\footnotesize\sffamily, >={Stealth[length=1.8mm]},
  stp/.style={draw, fill=ln-box, align=center, inner sep=3pt, minimum height=0.8cm, minimum width=4.4cm},
  miss/.style={stp, fill=white, minimum width=4.2cm},
  dec/.style={draw, diamond, aspect=2.4, fill=ln-accent!15, align=center, inner sep=1pt},
  lab/.style={font=\sffamily\scriptsize, text=ln-muted}]
  % --- hit path (left column)
  \node[stp] (split) at (0,0)
    {\iflangja{アドレスをタグ，インデックス，\\オフセットに分ける}{split the address into\\tag, index and offset}};
  \node[stp] (cmp) at (0,-1.3)
    {\iflangja{インデックスでセットを選び，\\有効な各ラインのタグと比較}{select the set by the index; compare\\the tag with every valid line of the set}};
  \node[dec] (hit) at (0,-2.65) {\iflangja{ヒット?}{hit?}};
  \node[stp] (acc) at (0,-6.55)
    {\iflangja{オフセットの位置を読み書き\\（書き込みなら D $\gets 1$）\\LRU カウンタを更新}{read or write at the offset\\(on a write: D $\gets 1$);\\update the LRU counters}};
  % --- miss path (right column)
  \node[miss] (victim) at (5.8,-2.65)
    {\iflangja{追い出すライン：無効なライン，\\なければ LRU カウンタが 0 のライン}{victim: an invalid line, or else\\the line whose LRU counter is 0}};
  \node[dec] (dirty) at (5.8,-3.95) {\iflangja{有効かつ\\ダーティ?}{valid and\\dirty?}};
  \node[miss] (wb) at (5.8,-5.25)
    {\iflangja{追い出すブロックを\\主記憶に書き戻す}{write the victim block\\back to main memory}};
  \node[miss] (load) at (5.8,-6.55)
    {\iflangja{ブロックを読み込みタグを格納\\V $\gets 1$，D $\gets 0$}{load the block, store its tag,\\V $\gets 1$, D $\gets 0$}};
  % --- arrows
  \draw[->] (split) -- (cmp);
  \draw[->] (cmp) -- (hit);
  \draw[->] (hit) -- node[lab, above] {\iflangja{いいえ}{no}} (victim);
  \draw[->] (hit) -- node[lab, right, pos=0.3] {\iflangja{はい}{yes}} (acc);
  \draw[->] (victim) -- (dirty);
  \draw[->] (dirty) -- node[lab, right] {\iflangja{はい}{yes}} (wb);
  \draw[->] (wb) -- (load);
  \draw[->] (dirty.east) -- node[lab, above] {\iflangja{いいえ}{no}} (8.4,-3.95) |- (load.east);
  \draw[->] (load) -- (acc);
\end{tikzpicture}

  \caption{LRU 置換とライトバック方式を用いるセットアソシアティブ
    キャッシュにおける1回のメモリアクセス．}
  \label{fig:l12-access-flow}
\end{figure}

\begin{example}
  \label{ex:l12-summary-trace}
  \cref{ex:l12-sa-conflict} の 2 ウェイセットアソシアティブキャッシュ
  （16 バイトのブロック，インデックスはビット 5--4，タグはビット 31--6）に，
  LRU 置換とライトバック方式を加えたものを考える．
  \cref{tab:l12-summary-trace} は，そのセット~1 を 5 回のアクセスにわたって
  追ったものである．セットの2つのラインは終始有効である．最初のアクセスは
  ヒットし，カウンタを更新するだけである．2番目のアクセスはミスし，
  カウンタが~0 のウェイ~1 が追い出される．そのブロックはダーティなので，
  新しいブロックを読み込む前に書き戻される．4番目のアクセスは書き込み
  ミスであり，追い出されるブロックはクリーンなので何も書き戻されず，
  読み込まれたブロックはただちにダーティになる．最後のアクセスは，3番目の
  アクセスが書き込んだタグ \texttt{0x40} のブロックを追い出し，それを
  書き戻す．
\end{example}

\begin{table}[htbp]
  \centering\small
  \caption{LRU 置換とライトバックを用いる 2 ウェイセットアソシアティブ
    キャッシュのセット~1（\cref{ex:l12-summary-trace}）．各アクセス後の
    各ウェイのタグ，ダーティビット，LRU カウンタを示す．メモリ操作のものも
    含め，タグはすべて 16 進数である．}
  \label{tab:l12-summary-trace}
  \setlength{\tabcolsep}{4pt}
  \begin{tabular}{@{}llccccccl@{}}
    \toprule
    & & \multicolumn{3}{c}{ウェイ 0} & \multicolumn{3}{c}{ウェイ 1} & \\
    \cmidrule(lr){3-5}\cmidrule(lr){6-8}
    アクセス & 結果 & タグ & D & LRU & タグ & D & LRU & メモリ操作 \\
    \midrule
    （初期状態）                &        & 40 & 0 & 0 & 43 & 1 & 1 & \\
    読み出し \texttt{0x1018}  & ヒット & 40 & 0 & 1 & 43 & 1 & 0 & --- \\
    読み出し \texttt{0x1054}  & ミス   & 40 & 0 & 0 & 41 & 0 & 1 & 43 書き戻し，41 読み込み \\
    書き込み \texttt{0x101C}  & ヒット & 40 & 1 & 1 & 41 & 0 & 0 & --- \\
    書き込み \texttt{0x1098}  & ミス   & 40 & 1 & 0 & 42 & 1 & 1 & 42 読み込み \\
    読み出し \texttt{0x1058}  & ミス   & 41 & 0 & 1 & 42 & 1 & 0 & 40 書き戻し，41 読み込み \\
    \bottomrule
  \end{tabular}
\end{table}

\begin{example}
  \label{ex:l12-overhead}
  \cref{tab:l12-organizations} は，128 バイトのブロック，32 ビットの
  アドレス，LRU 置換，ライトバックを用いる 64~KB のキャッシュの構成を，
  いくつかの連想度について求めたものである．8 ウェイのキャッシュでは，
  アドレス \texttt{0x0040A7C6}（ブロック番号 \texttt{0x814F}）の
  インデックスは $\texttt{0x814F} \mathbin{\&} \texttt{0x3F} = 15$，
  タグは $\texttt{0x814F} \gg 6 = \texttt{0x205}$ である．その 512 個の
  ラインには $512 \times 24 = \num{12288}$ ビット，すなわち 1536 バイトの
  メタデータが必要であり，これは 64~KB のデータの約 \SI{2.3}{\percent} に
  あたる．
\end{example}

\begin{table}[htbp]
  \centering\small
  \caption{128 バイトのブロックと 32 ビットのアドレスをもつ 64~KB の
    キャッシュのアドレスフィールドと，ライン当たりのメタデータのビット数
    （有効ビット $+$ ダーティビット $+$ LRU カウンタ $+$ タグ）．}
  \label{tab:l12-organizations}
  \begin{tabular}{@{}lrcccc@{}}
    \toprule
    構成 & セット数 & オフセット & インデックス & タグ & ライン当たりのメタデータ \\
    \midrule
    ダイレクトマップ     & 512 & 7 & 9 & 16 & $1+1+0+16 = 18$ \\
    2 ウェイ             & 256 & 7 & 8 & 17 & $1+1+1+17 = 20$ \\
    8 ウェイ             & 64  & 7 & 6 & 19 & $1+1+3+19 = 24$ \\
    フルアソシアティブ   & 1   & 7 & 0 & 25 & $1+1+9+25 = 36$ \\
    \bottomrule
  \end{tabular}
\end{table}

第13章では同じ考え方---ブロック，タグ，置換---を仮想アドレスの変換に
適用し，第14章では，ここで述べた基本的なキャッシュのヒット時間，ミス率，
ミスペナルティを下げる技法を扱う．

\begin{keypoint}[章のまとめ]
  プログラムは時間的局所性と空間的局所性をもってメモリにアクセスするので，
  最近使われたブロックを保持する小容量で高速なキャッシュが大半のアクセスに
  応える．その効果は
  $\text{AMAT} = \text{ヒット時間} + \text{ミス率} \times \text{ミスペナルティ}$
  （\cref{eq:l12-amat}）で測られる．メモリは2のべき乗の大きさのブロックに
  分割されてキャッシュのラインに置かれる．アドレスはタグ，インデックス，
  オフセットに分かれ，ラインが有効でそのタグが一致すればヒットとなる．
  ダイレクトマップ，セットアソシアティブ，フルアソシアティブの各キャッシュ
  は，ブロックが占めうるラインの数が異なり，ヒット時間とミス率を
  トレードオフする．ミスの際には LRU が最も長い間使われていないブロックを
  追い出し，ダーティビットを用いるライトバック方式は，変更されたブロックが
  追い出されるときにだけブロックをメモリに書き込む．
\end{keypoint}

\end{document}
